\section{Data}
\label{sec:data}

Data were collected from multiple sources to construct a municipality-level dataset of judicial statistics, including the number of people sentenced to prison, ordered to pay fines, processed, indicted, and placed in pretrial detention, among others. This section provides a detailed description of each data source, along with the procedures followed to obtain the information. Table \ref{tab:desc_stats_all} provides descriptive statistics of all the variables in the analysis.

\textbf{Sentencing data} For sentencing information, I use the Criminal Court Proceedings dataset provided by INEGI, which covers records from criminal and mixed first-instance courts under state jurisdiction—and, in the case of the former Federal District, under federal jurisdiction—for the period 1997–2012. It is important to distinguish between two separate datasets within this source.

The pre-trial dataset contains information on all individuals formally processed in the judicial system. These individuals have been captured and are subject to an active investigation by prosecutors. They face a control judge who determines their judicial status, such as whether they are placed in pretrial detention (\textit{formal prision}), allowed to continue the process without detention (\textit{sujeción a proceso}), or released (\textit{no sujeción a proceso, libertad}). This dataset captures early judicial decisions and also reflects both, policing and prosecutorial activity.

The criminal courts dataset includes individuals who have undergone judicial proceedings and have received a final sentencing decision (trial). Sentences vary, including acquittals (release), convictions resulting in prison terms, or fines imposed after a guilty verdict. Both datasets provide rich individual-level details, such as race, religion, occupation, municipality of residence, sentencing length, sentencing date, and characteristics of the crimes involved.

For analysis, I aggregate data at the municipality level and focus primarily on key variables relevant to policing and judicial outcomes. To explore changes in policing behavior, the pre-trial hearings dataset offers the total number of individuals processed, which corresponds closely to the prosecutorial decisions made by the executive branch. This dataset also records initial judicial rulings, including whether individuals are sent to pretrial detention or allowed to remain free during the judicial process. Such decisions may be influenced by prison capacity constraints, making variables like the number of individuals sentenced to pretrial detention or released without detention particularly insightful.

Finally, other critical variables include the number of individuals sentenced to prison and the total number of people sentenced to any outcome—whether release, imprisonment, or fines. These variables serve as the primary measures to assess judicial responses and outcomes at the municipal level. 

% For the period following the 2012 judicial reform, I rely on a complementary dataset also provided by INEGI, titled Law Enforcement in Criminal Matters. This dataset draws on administrative records from the criminal first-instance courts of each state’s Superior Court of Justice. It is designed to complement the National Census of State Law Enforcement and provides data on criminal cases at various stages of the judicial process, including the number and characteristics of crimes, victims, defendants, and convicted individuals.

\textbf{Prison capacity and population} I constructed a novel bimonthly panel dataset of all federal and state prisons in Mexico using PDFs from the Ministry of Prevention and Social Reintegration (\textit{Secretaría de Prevención y Readaptación Social}). Approximately 10,000 PDF pages containing prison statistics were OCR-processed, followed by manual correction to address errors. This dataset spans from 2000 to 2025 and includes detailed information on the opening and closing of prison facilities, as well as time-varying statistics such as installed capacity, federal and state prison populations, and total inmate counts.

Over time, prisons may have changed names while remaining operational or been opened and closed at different points. To address this, I applied a deduplication process that combined string-matching algorithms with thorough manual verification. Additionally, I geolocated each prison facility with precise coordinates using a mix of manual searches and automated tools, including the Google Maps API and OpenStreetMap (OSM). This careful deduplication and geolocation were critical because judicial sentencing data are matched at the municipality level to the nearest operational prison, whether federal or state.

The treatment variable was constructed by combining this prison dataset with information on individuals’ municipalities of residence. After geolocating each federal prison, I calculated the distance from each of Mexico’s 2,640 municipalities to all federal prisons over time, identifying the closest federal prison to each municipality in each period. Based on this, I created a time-varying dummy variable that equals one in all months following the opening of the nearest federal prison, provided the prison lies within a predetermined distance threshold.

It is important to note that, since the exact prison to which each individual is assigned is not observed, treatment assignment relies on a municipality-level prediction. This approach is justified by two key institutional facts: (1) individuals convicted of federal crimes are generally sent to federal prisons, and (2) the National Law on the Execution of Criminal Sentences mandates that sentenced individuals serve their time in the facility closest to their residence. These institutional rules support using geographic proximity as a valid proxy for prison assignment. Potential threats to identification related to this assumption will be discussed in the following section.

\textbf{Municipal characteristics}  Finally, to help with deviations from the parallel trends assumption and enhance the precision of my estimates, I use data from the National Census of Population and Housing 2000, conducted by INEGI every ten years. This census provides relevant information for each municipality, including the population, the distribution of urban and rural households, average education levels, and more.

